\documentclass{beamer}
\usepackage{hyperref}
\usepackage{amsmath}
\usepackage{amsfonts}
\usepackage{amssymb}
\title{Instructor Notes: Introduction to Global Shifts and Northern British Columbia}

\subtitle{Economics of Global Shifts: ECON 220}
\author{Muhebullah Karimzada}
\date{Winter 2025}

\begin{document}

\maketitle

\begin{frame}{Overview of the Lecture}
    \textbf{Lecture Focus:} 
    \begin{itemize}
        \item Understanding the concept of global economic shifts.
        \item Analyzing the impact of global shifts on Northern British Columbia (NBC).
        \item Discussing how NBC can leverage its resources and address local/global challenges.
    \end{itemize}
    \textbf{Instructor’s Goal:}
    \begin{itemize}
        \item Explain the causes and effects of global economic shifts.
        \item Link global shifts to local realities in NBC.
        \item Provide examples of how NBC can respond to global economic trends.
    \end{itemize}
\end{frame}

% Slide 2: Global Economic Shifts – Key Concept
\begin{frame}{Global Economic Shifts}
    \textbf{Definition:}  
    \begin{itemize}
        \item Global economic shifts are significant and sustained changes in the global economic environment that affect production, trade, labor markets, and social structures.
        \item These shifts can arise from technological advancements, policy changes, environmental factors, and demographic trends.
    \end{itemize}
    \textbf{Instructor's Example:}
    \begin{itemize}
        \item The Industrial Revolution: Shift from agrarian economies to industrial economies had lasting impacts on global trade patterns and labor markets.
    \end{itemize}
\end{frame}

% Slide 3: Why Global Shifts Matter
\begin{frame}{Why Global Shifts Matter}
    \begin{itemize}
        \item Impact global trade flows and international relations.
        \item Affect national economies, industries, and employment sectors.
        \item Influence regional economies by changing demand for resources, altering labor markets, and influencing government policies.
    \end{itemize}
    \textbf{Instructor's Example:}
    \begin{itemize}
        \item 2008 Financial Crisis: Disrupted global markets, leading to changes in housing prices, unemployment, and trade flows.
    \end{itemize}
\end{frame}

% Slide 4: Key Drivers of Global Economic Shifts
\begin{frame}{Drivers of Global Economic Shifts}
    \begin{itemize}
        \item \textbf{Technological Advancements:} Automation, AI, IoT, 3D printing.
        \item \textbf{Trade and Economic Integration:} Trade agreements (e.g., NAFTA, CPTPP).
        \item \textbf{Environmental and Climate Changes:} Sustainability goals, clean technology.
        \item \textbf{Demographic Changes:} Aging populations, urbanization.
        \item \textbf{Geopolitical Dynamics:} Trade wars, economic sanctions.
    \end{itemize}
    \textbf{Instructor’s Goal:}
    \begin{itemize}
        \item Explain each driver with relevant examples.
        \item Discuss how these drivers impact NBC’s economy, particularly in the context of natural resource extraction.
    \end{itemize}
\end{frame}

% Slide 5: Technological Advancements
\begin{frame}{Technological Advancements}
    \textbf{Impact of Technology on NBC:}
    \begin{itemize}
        \item Automation in forestry can improve efficiency but displace jobs in rural areas.
        \item AI and IoT could lead to new industries such as tech and innovation hubs.
    \end{itemize}
    \textbf{Instructor’s Example:}
    \begin{itemize}
        \item Automated machinery in forestry increases productivity but may reduce local job opportunities.
        \item How does automation impact labor markets in resource-dependent regions like NBC?
    \end{itemize}
\end{frame}

% Slide 6: Trade and Economic Integration
\begin{frame}{Trade and Economic Integration}
    \begin{itemize}
        \item Trade agreements such as CPTPP and NAFTA expand market access.
        \item NBC benefits from its proximity to the Asia-Pacific region, especially for LNG exports.
    \end{itemize}
    \textbf{Instructor’s Example:}
    \begin{itemize}
        \item The potential impact of the U.S.-China trade war on NBC’s export sectors (e.g., LNG, forestry).
        \item How can NBC diversify its markets to reduce reliance on a few trading partners?
    \end{itemize}
\end{frame}

% Slide 7: Environmental and Climate Changes
\begin{frame}{Environmental and Climate Changes}
    \begin{itemize}
        \item Global sustainability trends are reshaping industries like forestry, mining, and energy production.
        \item Environmental regulations add costs but can also drive innovation in green technologies.
    \end{itemize}
    \textbf{Instructor’s Example:}
    \begin{itemize}
        \item Hydroelectric power in NBC is a response to global sustainability trends.
        \item How can NBC balance resource extraction with environmental protection?
    \end{itemize}
\end{frame}

% Slide 8: Demographic Changes
\begin{frame}{Demographic Changes}
    \begin{itemize}
        \item Aging populations in developed countries, including Canada, increase demand for healthcare services.
        \item Outmigration of youth from NBC challenges workforce sustainability.
    \end{itemize}
    \textbf{Instructor’s Example:}
    \begin{itemize}
        \item Aging populations in urban centers create demand for healthcare, but outmigration from NBC leaves gaps in the workforce.
        \item How can NBC retain its youth and attract skilled labor?
    \end{itemize}
\end{frame}

% Slide 9: Geopolitical Dynamics
\begin{frame}{Geopolitical Dynamics}
    \begin{itemize}
        \item Regional conflicts, trade wars, and economic sanctions disrupt global trade and influence market access.
        \item NBC’s dependence on international markets makes it vulnerable to geopolitical changes.
    \end{itemize}
    \textbf{Instructor’s Example:}
    \begin{itemize}
        \item Trade tensions between the U.S. and China have affected demand for LNG from NBC.
        \item What can NBC do to diversify and reduce reliance on vulnerable markets?
    \end{itemize}
\end{frame}

% Slide 10: Northern British Columbia's Economic Profile
\begin{frame}{Northern British Columbia’s Economic Profile}
    \begin{itemize}
        \item Rich in natural resources: timber, natural gas, minerals.
        \item Access to the Asia-Pacific trade route via the Prince Rupert Port.
        \item Challenges: Vulnerability to commodity price volatility, aging population, and outmigration.
    \end{itemize}
    \textbf{Instructor’s Goal:}
    \begin{itemize}
        \item Explain the region’s strengths in natural resources and trade.
        \item Discuss challenges like aging demographics and reliance on resource extraction.
    \end{itemize}
\end{frame}

% Slide 11: Case Study - Prince Rupert Port
\begin{frame}{Case Study: Prince Rupert Port}
    \textbf{Background:}
    \begin{itemize}
        \item Canada’s fastest-growing port, connecting to Asia-Pacific markets.
        \item Strategic location for growing trade with China, Japan, and South Korea.
    \end{itemize}
    \textbf{Opportunities:}
    \begin{itemize}
        \item Expansion of port facilities increases trade volume.
        \item Economic growth through increased export capacity, especially for LNG.
    \end{itemize}
    \textbf{Instructor's Discussion:}
    \begin{itemize}
        \item How can NBC capitalize on its port to expand exports while addressing environmental and sustainability concerns?
    \end{itemize}
\end{frame}

% Slide 12: Indigenous Perspectives
\begin{frame}{Indigenous Perspectives}
    \begin{itemize}
        \item Indigenous communities are key stakeholders in NBC’s resource development.
        \item Indigenous-led projects, like joint ventures in LNG, provide opportunities for economic empowerment.
    \end{itemize}
    \textbf{Instructor’s Example:}
    \begin{itemize}
        \item Wet’suwet’en protests over pipeline construction highlight the complexities of balancing development with Indigenous rights.
    \end{itemize}
    \textbf{Instructor's Discussion:}
    \begin{itemize}
        \item How can NBC incorporate Indigenous perspectives in economic development for long-term sustainability?
    \end{itemize}
\end{frame}

% Slide 13: Class Discussion Questions
\begin{frame}{Class Discussion Questions}
    \begin{enumerate}
        \item What are the key global trends impacting Northern British Columbia?
        \item How can NBC balance economic growth with sustainability?
        \item What role should Indigenous communities play in NBC’s economic development?
    \end{enumerate}
\end{frame}

% Slide 14: Answer to Discussion Question 1
\begin{frame}{Answer to Discussion Question 1}
    \textbf{Key Global Trends Impacting NBC:}
    \begin{itemize}
        \item Global trade flows: Trade agreements and global supply chains impact export sectors.
        \item Technological advancements: Automation and green technologies shape the resource extraction and manufacturing industries.
        \item Geopolitical shifts: Trade tensions, sanctions, and geopolitical conflict can disrupt market access and impact exports.
    \end{itemize}
    \textbf{Instructor's Explanation:}
    \begin{itemize}
        \item NBC’s economy is highly dependent on natural resources like LNG and forestry. Global shifts in trade and technological innovation will impact demand and competition in these industries.
    \end{itemize}
\end{frame}

% Slide 15: Answer to Discussion Question 2
\begin{frame}{Answer to Discussion Question 2}
    \textbf{Balancing Economic Growth with Sustainability:}
    \begin{itemize}
        \item Prioritize resource efficiency: Implement green technologies, sustainable mining practices, and eco-friendly policies.
        \item Develop new industries: Focus on diversifying into clean energy, technology, and sustainable tourism.
        \item Collaborate with local communities: Ensure that development does not harm Indigenous rights and the environment.
    \end{itemize}
    \textbf{Instructor's Explanation:}
    \begin{itemize}
        \item Sustainable development requires balancing short-term economic gains with long-term environmental protection and social equity.
    \end{itemize}
\end{frame}

% Slide 16: Answer to Discussion Question 3
\begin{frame}{Answer to Discussion Question 3}
    \textbf{Role of Indigenous Communities in NBC's Economic Development:}
    \begin{itemize}
        \item Indigenous communities should be key partners in resource development projects.
        \item Collaborative models (e.g., joint ventures) empower Indigenous peoples, creating jobs and economic opportunities while respecting their rights.
        \item Indigenous knowledge and leadership are essential for sustainable development, especially in environmental conservation.
    \end{itemize}
    \textbf{Instructor's Explanation:}
    \begin{itemize}
        \item Incorporating Indigenous perspectives ensures that economic development respects cultural values and leads to long-term prosperity.
    \end{itemize}
\end{frame}

% Slide 17: Readings
\begin{frame}{Readings}
    \begin{itemize}
        \item Stiglitz, J. E. (2002). \textit{Globalization and Its Discontents}.
        \item Bowles, P., \& Wilson, M. (2015). \textit{Resource Communities in a Globalizing Region}.
        \item Prince Rupert Port Authority (2023). \textit{Annual Reports}.
    \end{itemize}
\end{frame}

\end{document}




